\chapter{Introdução aos sistemas de tipos}

\section{Introdução}

Nos capítulos anteriores, definimos semânticas operacionais para
linguagens de expressões e para linguagens imperativas. Vimos que a
semântica de uma linguagem pode ser dada por regras de inferência que
descrevem como os programas são executados. No entanto, a semântica
por si só não impede que expressões sem sentido sejam escritas: como
somar um booleano a um inteiro, ou executar
\haskell{pred true}. Essas expressões são
\textbf{sintaticamente válidas}, mas \textbf{semanticamente sem significado}.

É exatamente para isso que servem os \textbf{sistemas de tipos}:
impor, de forma estática, que apenas expressões semanticamente
bem-definidas possam ser executadas. Este capítulo apresenta as
diferentes estratégias de tipagem utilizadas por linguagens de
programação e discute os sistemas de tipos, utilizados para descrever
formalmente regras para validação de tipos em linguagens de programação.

\section{Tipos em Linguagens de Programação}

Um \textbf{tipo} é uma classificação de valores que define:

\begin{enumerate}
  \item Quais operações são aplicáveis a eles;
  \item Como eles são representados na memória;
  \item Que conjunto de valores o tipo compreende.
\end{enumerate}

Por exemplo, o tipo \haskell{Int} classifica valores
inteiros e admite operações aritméticas; o tipo
\haskell{Bool} classifica os valores
\haskell{true} e \haskell{false} e
admite operações lógicas. O \textbf{sistema de tipos} de uma
linguagem é o conjunto de regras que associa tipos a expressões e
verifica se essas associações são coerentes.

\subsection{Tipagem Estática e Dinâmica}

A distinção mais fundamental em sistemas de tipos é \emph{quando} a
verificação de tipos ocorre:

\textbf{Tipagem estática} (\emph{static typing}): a verificação é
feita em \textbf{tempo de compilação}, antes da execução do programa.
Se o programa contiver um erro de tipo, o compilador o rejeita. As
linguagens Haskell, Java, C e Rust são exemplos que usam tipagem estática.

\textbf{Tipagem dinâmica} (\emph{dynamic typing}): a verificação é
feita em \textbf{tempo de execução}, à medida que o programa é
executado. Erros de tipo são detectados apenas quando a operação
inválida é executada. Exemplos de linguagens que usam essa estratégia
são Python, JavaScript, Ruby e Lisp.

A tipagem estática oferece garantias mais fortes: o sistema de tipos
pode certificar que certos tipos de erros \emph{nunca} ocorrerão em
nenhuma execução do programa. A tipagem dinâmica oferece mais
flexibilidade, mas adia a detecção de erros para o momento da execução.

\subsection{Tipagem Forte e Fraca}

Uma distinção ortogonal é a de \textbf{tipagem forte} vs \textbf{tipagem fraca}:

\textbf{Tipagem forte} (\emph{strong typing}): a linguagem não
permite conversões implícitas entre tipos incompatíveis. Operações
entre tipos diferentes resultam em erro (em tempo de compilação ou
execução). Exemplos: Haskell, Python.

\textbf{Tipagem fraca} (\emph{weak typing}): a linguagem realiza
conversões implícitas entre tipos, frequentemente de maneiras
surpreendentes. Exemplos: C (converte implicitamente entre inteiros e
ponteiros), JavaScript (\haskell{1 + "1"} resulta em
\haskell{"11"}).

Note que estas classificações são independentes: Python é
dinamicamente tipado e fortemente tipado; C é estaticamente tipado e
fracamente tipado.

\section{A Linguagem de Expressões Tipadas}

Para discutirmos sobre sistemas de tipos, vamos utilizar um linguagem
de expressões simples. A sua \textbf{sintaxe abstrata} é representada
pela seguinte gramática:

\[
  \begin{array}{lcll}
    t & ::= & \mathbf{true}  & \text{(verdadeiro)} \\
    & \mid & \mathbf{false} & \text{(falso)} \\
    & \mid & \mathbf{if}\ t_1\ \mathbf{then}\ t_2\ \mathbf{else}\ t_3
    & \text{(condicional)} \\
    & \mid & \mathbf{0} & \text{(zero)} \\
    & \mid & \mathbf{succ}\ t & \text{(sucessor)} \\
    & \mid & \mathbf{pred}\ t & \text{(predecessor)} \\
    & \mid & \mathbf{iszero}\ t & \text{(teste de zero)}
  \end{array}
\]

Os \textbf{tipos} são:

\[
  \begin{array}{lcll}
    T & ::= & \mathbf{Bool} & \text{(booleano)} \\
    & \mid & \mathbf{Nat}  & \text{(número natural)}
  \end{array}
\]

Os \textbf{valores} da linguagem são os termos que já não podem ser
reduzidos pela semântica operacional. Distinguimos dois tipos de valores:

\[
  \begin{array}{lcll}
    v & ::= & \mathbf{true} \mid \mathbf{false} \mid nv \\
    nv & ::= & \mathbf{0} \mid \mathbf{succ}\ nv
  \end{array}
\]

onde \(nv\) denota os \textbf{valores numéricos}, isto é, \(\mathbf{0}\) e
qualquer número de aplicações de \(\mathbf{succ}\) a \(\mathbf{0}\).

\section{O Sistema de Tipos}

Um \textbf{sistema de tipos} é definido por uma \textbf{relação de
tipagem} entre termos e tipos. Escrevemos \(\vdash t : T\) para dizer
que ``o termo \(t\) tem tipo \(T\)''. A relação é definida por um
conjunto de regras de inferência.

\subsection{Regras de Tipagem}

\textbf{Constantes booleanas:}

\[
  \frac{}{\vdash \mathbf{true} : \mathbf{Bool}} \quad\text{(T-True)}
  \qquad
  \frac{}{\vdash \mathbf{false} : \mathbf{Bool}} \quad\text{(T-False)}
\]

\textbf{Condicional:}

\[
  \frac{
    \vdash t_1 : \mathbf{Bool} \quad
    \vdash t_2 : T \quad
    \vdash t_3 : T
  }{
    \vdash \mathbf{if}\ t_1\ \mathbf{then}\ t_2\ \mathbf{else}\ t_3 : T
  } \quad\text{(T-If)}
\]

A regra (T-If) exige que a condição \(t_1\) seja booleana e que os
dois ramos \(t_2\) e \(t_3\) tenham o \textbf{mesmo tipo} \(T\). O
resultado do condicional tem esse tipo.

\textbf{Naturais:}

\[
  \frac{}{\vdash \mathbf{0} : \mathbf{Nat}} \quad\text{(T-Zero)}
\]

\[
  \frac{\vdash t_1 : \mathbf{Nat}}{\vdash \mathbf{succ}\ t_1 :
  \mathbf{Nat}} \quad\text{(T-Succ)}
  \qquad
  \frac{\vdash t_1 : \mathbf{Nat}}{\vdash \mathbf{pred}\ t_1 :
  \mathbf{Nat}} \quad\text{(T-Pred)}
\]

\[
  \frac{\vdash t_1 : \mathbf{Nat}}{\vdash \mathbf{iszero}\ t_1 :
  \mathbf{Bool}} \quad\text{(T-IsZero)}
\]

\subsection{Exemplos de Derivações de Tipo}

\textbf{Exemplo 1:} O termo
\(\mathbf{iszero}\ (\mathbf{succ}\ \mathbf{0})\) tem tipo \(\mathbf{Bool}\):

\[
  \dfrac{
    \dfrac{
      \dfrac{}{\vdash \mathbf{0} : \mathbf{Nat}}
    }{\vdash \mathbf{succ}\ \mathbf{0} : \mathbf{Nat}}
  }{\vdash \mathbf{iszero}\ (\mathbf{succ}\ \mathbf{0}) : \mathbf{Bool}}
\]

\textbf{Exemplo 2:} O condicional
\(\mathbf{if}\ (\mathbf{iszero}\ \mathbf{0})\ \mathbf{then}\ \mathbf{0}\ \mathbf{else}\ \mathbf{succ}\ \mathbf{0}\)
tem tipo \(\mathbf{Nat}\):

\[
  \dfrac{
    \dfrac{
      \dfrac{}{\vdash \mathbf{0} : \mathbf{Nat}}
    }{\vdash \mathbf{iszero}\ \mathbf{0} : \mathbf{Bool}}
    \quad
    \dfrac{}{\vdash \mathbf{0} : \mathbf{Nat}}
    \quad
    \dfrac{
      \dfrac{}{\vdash \mathbf{0} : \mathbf{Nat}}
    }{\vdash \mathbf{succ}\ \mathbf{0} : \mathbf{Nat}}
  }{\vdash
    \mathbf{if}\ (\mathbf{iszero}\ \mathbf{0})\ \mathbf{then}\ \mathbf{0}\ \mathbf{else}\ \mathbf{succ}\ \mathbf{0}
  : \mathbf{Nat}}
\]

\textbf{Exemplo 3 (erro de tipo):} O termo
\(\mathbf{succ}\ \mathbf{true}\) \textbf{não é tipável}: a regra
(T-Succ) exige que o argumento tenha tipo \(\mathbf{Nat}\), mas
\(\mathbf{true}\) tem tipo \(\mathbf{Bool}\). Não existe derivação
válida para este termo.

\subsection{Unicidade de Tipos}

Uma propriedade importante desta linguagem é a \textbf{unicidade de
tipos}: cada termo bem tipado tem \textbf{exatamente um tipo}. Formalmente:

\begin{quote}
  \textbf{Lema (Unicidade):} Se \(\vdash t : T\) e \(\vdash t : T'\),
  então \(T = T'\).
\end{quote}

\textbf{Prova:} Por indução estrutural no termo \(t\).

\begin{itemize}
  \item
    \emph{Caso} \(t = \mathbf{true}\): a única regra aplicável é
    (T-True), que conclui \(T = \mathbf{Bool}\). Logo \(T = T' =
    \mathbf{Bool}\).
  \item
    \emph{Caso} \(t = \mathbf{false}\): análogo.
  \item
    \emph{Caso} \(t = \mathbf{0}\): a única regra aplicável é
    (T-Zero), que conclui \(T = \mathbf{Nat}\). Logo \(T = T' = \mathbf{Nat}\).
  \item
    \emph{Caso} \(t = \mathbf{succ}\ t_1\): a única regra aplicável é
    (T-Succ). Ambas as derivações de \(T\) e de \(T'\) devem usar
    (T-Succ), portanto \(T = T' = \mathbf{Nat}\).
  \item
    \emph{Caso} \(t = \mathbf{pred}\ t_1\): análogo; a única regra
    é (T-Pred), logo \(T = T' = \mathbf{Nat}\).
  \item
    \emph{Caso} \(t = \mathbf{iszero}\ t_1\): a única regra é
    (T-IsZero), logo \(T = T' = \mathbf{Bool}\).
  \item
    \emph{Caso} \(t =
    \mathbf{if}\ t_1\ \mathbf{then}\ t_2\ \mathbf{else}\ t_3\): a
    única regra aplicável é (T-If). Pela hipótese de indução aplicada
    a \(t_2\), o tipo de \(t_2\) é único; como ambas as derivações
    concluem \(T\) e \(T'\) com base no tipo de \(t_2\), temos \(T = T'\).
\end{itemize}

Como a prova cobre todos os casos da gramática, o lema está
demonstrado. A consequência prática é que a tipagem é
\textbf{determinística}: inspecionando um termo, é possível
determinar seu tipo de forma única, sem ambiguidade.

\section{Semântica Small-Step}

A semântica small-step define a execução como uma sequência de
reduções \(t \to t'\). Os termos que não podem ser reduzidos são os
\textbf{valores} \(v\) e as \textbf{formas presas} (\emph{stuck
terms}), termos sem mais reduções disponíveis que não são valores,
como \(\mathbf{succ}\ \mathbf{true}\).

\subsection{Regras de Redução para Booleanos}

\textbf{Condicional (redução da condição):}

\[
  \frac{t_1 \to
  t_1'}{\mathbf{if}\ t_1\ \mathbf{then}\ t_2\ \mathbf{else}\ t_3 \to
  \mathbf{if}\ t_1'\ \mathbf{then}\ t_2\ \mathbf{else}\ t_3}
  \quad\text{(E-If)}
\]

\textbf{Condicional (ramos):}

\[
  \frac{}{\mathbf{if}\ \mathbf{true}\ \mathbf{then}\ t_2\ \mathbf{else}\ t_3
  \to t_2}
  \quad\text{(E-IfTrue)}
  \qquad
  \frac{}{\mathbf{if}\ \mathbf{false}\ \mathbf{then}\ t_2\ \mathbf{else}\ t_3
  \to t_3}
  \quad\text{(E-IfFalse)}
\]

\subsection{Regras de Redução para Naturais}

\textbf{Sucessor:}

\[
  \frac{t_1 \to t_1'}{\mathbf{succ}\ t_1 \to \mathbf{succ}\ t_1'}
  \quad\text{(E-Succ)}
\]

\textbf{Predecessor:}

\[
  \frac{t_1 \to t_1'}{\mathbf{pred}\ t_1 \to \mathbf{pred}\ t_1'}
  \quad\text{(E-Pred)}
  \qquad
  \frac{}{\mathbf{pred}\ \mathbf{0} \to \mathbf{0}}
  \quad\text{(E-PredZero)}
  \qquad
  \frac{}{\mathbf{pred}\ (\mathbf{succ}\ nv_1) \to nv_1}
  \quad\text{(E-PredSucc)}
\]

\textbf{Teste de zero:}

\[
  \frac{t_1 \to t_1'}{\mathbf{iszero}\ t_1 \to \mathbf{iszero}\ t_1'}
  \quad\text{(E-IsZero)}
  \qquad
  \frac{}{\mathbf{iszero}\ \mathbf{0} \to \mathbf{true}}
  \quad\text{(E-IsZeroZero)}
  \qquad
  \frac{}{\mathbf{iszero}\ (\mathbf{succ}\ nv_1) \to \mathbf{false}}
  \quad\text{(E-IsZeroSucc)}
\]

\subsection{Formas Presas}

Uma \textbf{forma presa} (\emph{stuck term}) é um termo que não é um
valor e para o qual nenhuma regra de redução se aplica. Por exemplo:

\begin{itemize}
  \item \(\mathbf{succ}\ \mathbf{true}\): não é um valor (pois
    \(\mathbf{true}\) não é um \(nv\)) e nenhuma regra se aplica;
  \item \(\mathbf{pred}\ \mathbf{false}\): idem;
  \item \(\mathbf{if}\ \mathbf{0}\ \mathbf{then}\ t_2\ \mathbf{else}\ t_3\):
    \(\mathbf{0}\) é um valor numérico, não booleano.
\end{itemize}

Formas presas representam \textbf{erros em tempo de execução}. O
papel do sistema de tipos é garantir que termos bem tipados nunca se
tornem formas presas.

\subsection{Exemplo: Small-Step}

Considere o termo \(t =
\mathbf{pred}\ (\mathbf{succ}\ (\mathbf{succ}\ \mathbf{0}))\):

\[
  \mathbf{pred}\ (\mathbf{succ}\ (\mathbf{succ}\ \mathbf{0}))
  \xrightarrow{\text{E-PredSucc}}
  \mathbf{succ}\ \mathbf{0}
\]

O resultado \(\mathbf{succ}\ \mathbf{0}\) é um valor numérico (\(nv\)).

Considere agora:

\[
  \mathbf{if}\ (\mathbf{iszero}\ \mathbf{0})\ \mathbf{then}\ (\mathbf{succ}\ \mathbf{0})\ \mathbf{else}\ \mathbf{0}
\]

\[
  \xrightarrow{\text{E-If, E-IsZeroZero}}
  \mathbf{if}\ \mathbf{true}\ \mathbf{then}\ (\mathbf{succ}\ \mathbf{0})\ \mathbf{else}\ \mathbf{0}
\]

\[
  \xrightarrow{\text{E-IfTrue}}
  \mathbf{succ}\ \mathbf{0}
\]

\section{Semântica Big-Step}

A semântica big-step define a relação \(t \Downarrow v\), lida como
``o termo \(t\) avalia para o valor \(v\)''.

\subsection{Regras para Booleanos}

\[
  \frac{}{\mathbf{true} \Downarrow \mathbf{true}} \quad\text{(B-True)}
  \qquad
  \frac{}{\mathbf{false} \Downarrow \mathbf{false}} \quad\text{(B-False)}
\]

\[
  \frac{t_1 \Downarrow \mathbf{true} \quad t_2 \Downarrow v}
  {\mathbf{if}\ t_1\ \mathbf{then}\ t_2\ \mathbf{else}\ t_3 \Downarrow v}
  \quad\text{(B-IfTrue)}
  \qquad
  \frac{t_1 \Downarrow \mathbf{false} \quad t_3 \Downarrow v}
  {\mathbf{if}\ t_1\ \mathbf{then}\ t_2\ \mathbf{else}\ t_3 \Downarrow v}
  \quad\text{(B-IfFalse)}
\]

\subsection{Regras para Naturais}

\[
  \frac{}{\mathbf{0} \Downarrow \mathbf{0}} \quad\text{(B-Zero)}
  \qquad
  \frac{t_1 \Downarrow nv_1}{\mathbf{succ}\ t_1 \Downarrow
  \mathbf{succ}\ nv_1} \quad\text{(B-Succ)}
\]

\[
  \frac{t_1 \Downarrow \mathbf{0}}{\mathbf{pred}\ t_1 \Downarrow
  \mathbf{0}} \quad\text{(B-PredZero)}
  \qquad
  \frac{t_1 \Downarrow \mathbf{succ}\ nv_1}{\mathbf{pred}\ t_1
  \Downarrow nv_1} \quad\text{(B-PredSucc)}
\]

\[
  \frac{t_1 \Downarrow \mathbf{0}}{\mathbf{iszero}\ t_1 \Downarrow
  \mathbf{true}} \quad\text{(B-IsZeroZero)}
  \qquad
  \frac{t_1 \Downarrow \mathbf{succ}\ nv_1}{\mathbf{iszero}\ t_1
  \Downarrow \mathbf{false}} \quad\text{(B-IsZeroSucc)}
\]

\subsection{Exemplo: Big-Step}

\textbf{Exemplo:}
\(\mathbf{if}\ (\mathbf{iszero}\ \mathbf{0})\ \mathbf{then}\ \mathbf{succ}\ \mathbf{0}\ \mathbf{else}\ \mathbf{0}
\Downarrow \mathbf{succ}\ \mathbf{0}\)

\[
  \dfrac{
    \dfrac{
      \dfrac{}{\mathbf{0} \Downarrow \mathbf{0}}
    }{\mathbf{iszero}\ \mathbf{0} \Downarrow \mathbf{true}}
    \quad
    \dfrac{
      \dfrac{}{\mathbf{0} \Downarrow \mathbf{0}}
    }{\mathbf{succ}\ \mathbf{0} \Downarrow \mathbf{succ}\ \mathbf{0}}
  }{
    \mathbf{if}\ (\mathbf{iszero}\ \mathbf{0})\ \mathbf{then}\ \mathbf{succ}\ \mathbf{0}\ \mathbf{else}\ \mathbf{0}
    \Downarrow \mathbf{succ}\ \mathbf{0}
  }
\]

\section{Corretude do Sistema de Tipos}

A propriedade de \textbf{corretude} (\emph{type soundness})
estabelece que o sistema de tipos é uma garantia válida: termos bem
tipados nunca se tornam formas presas. Em outras palavras, um
programa que passa na verificação de tipos \textbf{sempre progride}
até um valor, sem travar em uma operação sem sentido.

Formalmente, a corretude é enunciada como:

\begin{quote}
  \textbf{Teorema (Type Soundness):} Se \(\vdash t : T\) e \(t \to^*
  t'\), então \(t'\) é um valor ou existe \(t''\) tal que \(t' \to t''\).
\end{quote}

Em português: \emph{programas bem tipados não ficam presos}. A prova
deste teorema se apoia em dois lemas: o de preservação e o de progresso.

\subsection{Lema de Preservação}

\begin{quote}
  \textbf{Lema de Preservação (\emph{Preservation}):} Se \(\vdash t : T\) e
  \(t \to t'\), então \(\vdash t' : T\).
\end{quote}

Este lema afirma que \textbf{a tipagem é invariante sob redução}: se
um termo tem tipo \(T\) e dá um passo de execução, o termo resultante
ainda tem o mesmo tipo \(T\). O tipo é \emph{preservado} ao longo de
toda a computação.

\textbf{Por que é importante?} Sem a preservação, um termo poderia
começar com tipo correto e, após algumas reduções, tornar-se um termo
de tipo diferente, o que quebraria qualquer garantia que o
verificador de tipos tentasse fornecer. A preservação garante que a
``etiqueta de tipo'' permanece coerente em cada passo.

\textbf{Prova:} Por indução na derivação de \(t \to t'\), analisando
cada regra de redução.

\emph{Caso} (E-IfTrue): \(t =
\mathbf{if}\ \mathbf{true}\ \mathbf{then}\ t_2\ \mathbf{else}\ t_3\)
e \(t' = t_2\). Como \(\vdash t : T\) foi derivado por (T-If), temos
\(\vdash \mathbf{true} : \mathbf{Bool}\), \(\vdash t_2 : T\) e
\(\vdash t_3 : T\). Logo \(\vdash t' = t_2 : T\).

\emph{Caso} (E-IfFalse): \(t =
\mathbf{if}\ \mathbf{false}\ \mathbf{then}\ t_2\ \mathbf{else}\ t_3\)
e \(t' = t_3\). Analogamente, de (T-If) obtemos \(\vdash t_3 : T\),
portanto \(\vdash t' : T\).

\emph{Caso} (E-If): \(t =
\mathbf{if}\ t_1\ \mathbf{then}\ t_2\ \mathbf{else}\ t_3\) com \(t_1
\to t_1'\) e \(t' =
\mathbf{if}\ t_1'\ \mathbf{then}\ t_2\ \mathbf{else}\ t_3\). De
(T-If) temos \(\vdash t_1 : \mathbf{Bool}\), \(\vdash t_2 : T\) e
\(\vdash t_3 : T\). Pela hipótese de indução aplicada a \(t_1 \to
t_1'\), obtemos \(\vdash t_1' : \mathbf{Bool}\). Aplicando (T-If),
concluímos \(\vdash t' : T\).

\emph{Caso} (E-Succ): \(t = \mathbf{succ}\ t_1\) com \(t_1 \to t_1'\)
e \(t' = \mathbf{succ}\ t_1'\). De (T-Succ) temos \(\vdash t_1 :
\mathbf{Nat}\). Pela hipótese de indução, \(\vdash t_1' :
\mathbf{Nat}\). Aplicando (T-Succ), concluímos \(\vdash
\mathbf{succ}\ t_1' : \mathbf{Nat}\).

\emph{Caso} (E-Pred): \(t = \mathbf{pred}\ t_1\) com \(t_1 \to t_1'\)
e \(t' = \mathbf{pred}\ t_1'\). De (T-Pred) temos \(\vdash t_1 :
\mathbf{Nat}\). Pela hipótese de indução, \(\vdash t_1' :
\mathbf{Nat}\). Aplicando (T-Pred), concluímos \(\vdash
\mathbf{pred}\ t_1' : \mathbf{Nat}\).

\emph{Caso} (E-PredZero): \(t = \mathbf{pred}\ \mathbf{0}\) e \(t' =
\mathbf{0}\). De (T-Pred) e (T-Zero) temos \(\vdash \mathbf{0} :
\mathbf{Nat}\), logo \(\vdash t' : \mathbf{Nat}\).

\emph{Caso} (E-PredSucc): \(t =
\mathbf{pred}\ (\mathbf{succ}\ nv_1)\) e \(t' = nv_1\). De (T-Pred)
temos \(\vdash \mathbf{succ}\ nv_1 : \mathbf{Nat}\); de (T-Succ)
obtemos \(\vdash nv_1 : \mathbf{Nat}\). Logo \(\vdash t' : \mathbf{Nat}\).

\emph{Caso} (E-IsZero): \(t = \mathbf{iszero}\ t_1\) com \(t_1 \to
t_1'\) e \(t' = \mathbf{iszero}\ t_1'\). De (T-IsZero) temos \(\vdash
t_1 : \mathbf{Nat}\). Pela hipótese de indução, \(\vdash t_1' :
\mathbf{Nat}\). Aplicando (T-IsZero), concluímos \(\vdash
\mathbf{iszero}\ t_1' : \mathbf{Bool}\).

\emph{Caso} (E-IsZeroZero): \(t = \mathbf{iszero}\ \mathbf{0}\) e
\(t' = \mathbf{true}\). Por (T-True), \(\vdash \mathbf{true} :
\mathbf{Bool}\); e de (T-IsZero) o tipo de \(t\) é \(\mathbf{Bool}\).
Logo \(\vdash t' : \mathbf{Bool}\).

\emph{Caso} (E-IsZeroSucc): \(t =
\mathbf{iszero}\ (\mathbf{succ}\ nv_1)\) e \(t' = \mathbf{false}\).
Por (T-False), \(\vdash \mathbf{false} : \mathbf{Bool}\); e de
(T-IsZero) o tipo de \(t\) é \(\mathbf{Bool}\). Logo \(\vdash t' :
\mathbf{Bool}\).

Como todos os casos da relação de redução estão cobertos, o lema está
demonstrado.

\subsection{Lema de Progresso}

\begin{quote}
  \textbf{Lema de Progresso (\emph{Progress}):} Se \(\vdash t : T\), então
  \(t\) é um valor ou existe \(t'\) tal que \(t \to t'\).
\end{quote}

Este lema afirma que \textbf{termos bem tipados nunca ficam presos}:
a cada passo de execução, ou o termo já é um valor (a computação
terminou), ou existe pelo menos uma redução possível (a computação
pode continuar).

\textbf{Por que é importante?} Um termo preso é um programa em um
estado de erro irrecuperável, como tentar aplicar
\haskell{pred} a um booleano. O lema de progresso
garante que isso nunca acontece para termos bem tipados: o sistema de
tipos elimina exatamente os termos que poderiam ficar presos.

\textbf{Prova:} Por indução estrutural na derivação de \(\vdash t : T\).

\emph{Caso} (T-True): \(t = \mathbf{true}\) é um valor.

\emph{Caso} (T-False): \(t = \mathbf{false}\) é um valor.

\emph{Caso} (T-Zero): \(t = \mathbf{0}\) é um valor numérico, portanto um valor.

\emph{Caso} (T-If): \(t =
\mathbf{if}\ t_1\ \mathbf{then}\ t_2\ \mathbf{else}\ t_3\) com
\(\vdash t_1 : \mathbf{Bool}\). Pela hipótese de indução aplicada a
\(t_1\): ou \(t_1\) é valor, ou existe \(t_1'\) tal que \(t_1 \to t_1'\).

\begin{itemize}
  \item Se \(t_1 \to t_1'\): aplica-se (E-If), obtendo \(t \to
    \mathbf{if}\ t_1'\ \mathbf{then}\ t_2\ \mathbf{else}\ t_3\).
  \item Se \(t_1\) é valor de tipo \(\mathbf{Bool}\): pelo lema de
    unicidade de tipos, \(t_1\) não pode ser um valor numérico \(nv\)
    (que tem tipo \(\mathbf{Nat}\)). Portanto \(t_1 \in
    \{\mathbf{true}, \mathbf{false}\}\).

    \begin{itemize}
      \item Se \(t_1 = \mathbf{true}\): aplica-se (E-IfTrue), \(t \to t_2\).
      \item Se \(t_1 = \mathbf{false}\): aplica-se (E-IfFalse), \(t \to t_3\).
    \end{itemize}
\end{itemize}

\emph{Caso} (T-Succ): \(t = \mathbf{succ}\ t_1\) com \(\vdash t_1 :
\mathbf{Nat}\). Pela hipótese de indução aplicada a \(t_1\): ou
\(t_1\) é valor, ou reduz.

\begin{itemize}
  \item Se \(t_1 \to t_1'\): aplica-se (E-Succ), \(t \to \mathbf{succ}\ t_1'\).
  \item Se \(t_1\) é valor de tipo \(\mathbf{Nat}\): pelo lema de
    unicidade, \(t_1\) não pode ser \(\mathbf{true}\) nem
    \(\mathbf{false}\) (que têm tipo \(\mathbf{Bool}\)), portanto
    \(t_1\) é um valor numérico \(nv\). Logo \(t =
    \mathbf{succ}\ nv\) é também um valor numérico.
\end{itemize}

\emph{Caso} (T-Pred): \(t = \mathbf{pred}\ t_1\) com \(\vdash t_1 :
\mathbf{Nat}\). Pela hipótese de indução aplicada a \(t_1\): ou
\(t_1\) é valor, ou reduz.

\begin{itemize}
  \item Se \(t_1 \to t_1'\): aplica-se (E-Pred), \(t \to \mathbf{pred}\ t_1'\).
  \item Se \(t_1\) é valor de tipo \(\mathbf{Nat}\): pelo lema de
    unicidade, \(t_1\) é um valor numérico \(nv\). Há dois subcasos:
    \begin{itemize}
      \item \(nv = \mathbf{0}\): aplica-se (E-PredZero), \(t \to \mathbf{0}\).
      \item \(nv = \mathbf{succ}\ nv_1\): aplica-se (E-PredSucc), \(t
        \to nv_1\).
    \end{itemize}
\end{itemize}

\emph{Caso} (T-IsZero): \(t = \mathbf{iszero}\ t_1\) com \(\vdash t_1
: \mathbf{Nat}\). Pela hipótese de indução aplicada a \(t_1\): ou
\(t_1\) é valor, ou reduz.
\begin{itemize}
  \item Se \(t_1 \to t_1'\): aplica-se (E-IsZero), \(t \to
    \mathbf{iszero}\ t_1'\).
  \item Se \(t_1\) é valor de tipo \(\mathbf{Nat}\): pelo lema de
    unicidade, \(t_1\) é um valor numérico \(nv\). Há dois subcasos:
    \begin{itemize}
      \item \(nv = \mathbf{0}\): aplica-se (E-IsZeroZero), \(t \to
        \mathbf{true}\).
      \item \(nv = \mathbf{succ}\ nv_1\): aplica-se (E-IsZeroSucc), \(t
        \to \mathbf{false}\).
    \end{itemize}
\end{itemize}

Como todos os casos da derivação de tipagem estão cobertos, o lema
está demonstrado.

\subsection{Da Preservação e do Progresso à Corretude}

Com os dois lemas, a prova de corretude é imediata:

\begin{quote}
  \textbf{Prova do Teorema (Type Soundness):} Seja \(\vdash t : T\) e
  \(t \to^* t'\). Queremos mostrar que \(t'\) é valor ou pode reduzir.

  Por \textbf{preservação}, como \(\vdash t : T\) e \(t \to^* t'\),
  temos \(\vdash t' : T\) (aplicando o lema de preservação em cada
  passo da sequência de reduções, por indução no comprimento da sequência).

  Por \textbf{progresso}, como \(\vdash t' : T\), \(t'\) é valor ou
  existe \(t''\) tal que \(t' \to t''\).
\end{quote}

A contribuição de cada lema é distinta e indispensável:

\begin{table}[H]
  \begin{tabular}{
      >{\raggedright\arraybackslash}p{0.20\linewidth}
      >{\raggedright\arraybackslash}p{0.35\linewidth}
      >{\raggedright\arraybackslash}p{0.35\linewidth}
    }
    \toprule
    \textbf{Lema} & \textbf{Garante} & \textbf{Sem ele, poderia acontecer} \\
    \midrule
    Preservação & O tipo se mantém durante a execução & Um passo de
    redução muda o tipo e o progresso não se aplica ao resultado \\
    \addlinespace
    Progresso & Termos bem tipados não ficam presos & O tipo se
    mantém, mas a execução trava mesmo assim \\
    \bottomrule
  \end{tabular}
  \centering
\end{table}

A vantagem desta abordagem é que ela \textbf{decompõe a corretude} em
dois lemas independentes e mais simples, cada um demonstrável por
indução. A prova de progresso usa \emph{apenas} as regras de tipagem;
a prova de preservação usa \emph{apenas} as regras de tipagem e as de redução.

% ============================================================
\section{Sistema de Tipos para TLine}
% ============================================================

TLine é uma linguagem imperativa simples com declaração de variáveis
tipadas, atribuição, leitura e impressão. Os seus tipos são
\(\mathbf{Int}\), \(\mathbf{Bool}\) e \(\mathbf{String}\). A
principal novidade em relação a TExp é a presença de
\textbf{variáveis}, o que torna necessário um \textbf{contexto de tipagem}:

\[
  \Gamma : \mathit{Var} \rightharpoonup \mathit{Ty}
\]

\subsection{Regras de Tipagem para Expressões}

As regras têm a forma \(\Gamma \vdash e : T\):

\[
  \frac{}{\Gamma \vdash n : \mathbf{Int}} \quad\text{(T-Int)}
  \qquad
  \frac{}{\Gamma \vdash b : \mathbf{Bool}} \quad\text{(T-Bool)}
  \qquad
  \frac{}{\Gamma \vdash s : \mathbf{String}} \quad\text{(T-String)}
\]

\[
  \frac{x \in \mathrm{dom}(\Gamma)}{\Gamma \vdash x : \Gamma(x)}
  \quad\text{(T-Var)}
\]

\[
  \frac{\Gamma \vdash e_1 : \mathbf{Int} \quad \Gamma \vdash e_2 : \mathbf{Int}}
  {\Gamma \vdash e_1 \oplus e_2 : \mathbf{Int}} \quad\text{(T-ArithOp)}
  \qquad
  \frac{\Gamma \vdash e_1 : \mathbf{Int} \quad \Gamma \vdash e_2 : \mathbf{Int}}
  {\Gamma \vdash e_1 \otimes e_2 : \mathbf{Bool}} \quad\text{(T-CmpOp)}
\]

\[
  \frac{\Gamma \vdash e : \mathbf{Bool}}
  {\Gamma \vdash \neg e : \mathbf{Bool}} \quad\text{(T-Not)}
\]

onde \(\oplus \in \{+, -, \star, /\}\) e \(\otimes \in \{<, =\}\).

\subsection{Regras de Tipagem para Comandos}

As regras têm a forma \(\Gamma \vdash s \dashv \Gamma'\), onde
\(\Gamma'\) é o contexto produzido após verificar o comando \(s\):

\[
  \frac{\Gamma \vdash e : T}
  {\Gamma \vdash \mathbf{var}\ x : T = e \dashv \Gamma[x \mapsto T]}
  \quad\text{(T-Decl)}
\]

\[
  \frac{x \in \mathrm{dom}(\Gamma) \quad \Gamma \vdash e : \Gamma(x)}
  {\Gamma \vdash x \mathbin{:=} e \dashv \Gamma} \quad\text{(T-Assign)}
\]

\[
  \frac{\Gamma \vdash e_p : \mathbf{String} \quad x \in
  \mathrm{dom}(\Gamma) \quad \Gamma(x) = \mathbf{Int}}
  {\Gamma \vdash \mathbf{read}\ e_p\ x \dashv \Gamma} \quad\text{(T-Read)}
\]

\[
  \frac{\Gamma \vdash e : T}
  {\Gamma \vdash \mathbf{print}\ e \dashv \Gamma} \quad\text{(T-Print)}
\]

A regra (T-Decl) é a única que \textbf{estende} o contexto: declara
\(x\) com tipo \(T\), produzindo \(\Gamma' = \Gamma[x \mapsto T]\).
As demais \textbf{preservam} o contexto: \(\Gamma' = \Gamma\).

% ============================================================
\section{Sistema de Tipos para TWhile}
% ============================================================

TWhile estende TLine com condicionais e repetição:

\begin{itemize}
  \item \(\mathbf{if}\ e\ \mathbf{then}\ B_1\ \mathbf{else}\ B_2\):
    executa o bloco \(B_1\) se \(e\) for verdadeiro, \(B_2\) caso contrário.
  \item \(\mathbf{while}\ e\ \mathbf{do}\ B\): repete o bloco \(B\)
    enquanto \(e\) for verdadeiro.
\end{itemize}

A novidade semântica fundamental é o \textbf{escopo léxico}: variáveis
declaradas dentro de um bloco \(\mathbf{if}\) ou \(\mathbf{while}\)
não são visíveis fora dele. Isso é formalizado pelas regras de tipagem:
o contexto de \textbf{saída} é sempre o contexto de
\textbf{entrada} \(\Gamma\), independentemente das declarações
internas aos blocos.

\subsection{Regras de Tipagem para os Novos Comandos}

\[
  \frac{\Gamma \vdash e : \mathbf{Bool}
    \quad \Gamma \vdash B_1 \dashv \Gamma_1
  \quad \Gamma \vdash B_2 \dashv \Gamma_2}
  {\Gamma \vdash
  \mathbf{if}\ e\ \mathbf{then}\ B_1\ \mathbf{else}\ B_2 \dashv \Gamma}
  \quad\text{(T-If)}
\]

\[
  \frac{\Gamma \vdash e : \mathbf{Bool}
  \quad \Gamma \vdash B \dashv \Gamma'}
  {\Gamma \vdash \mathbf{while}\ e\ \mathbf{do}\ B \dashv \Gamma}
  \quad\text{(T-While)}
\]

O contexto de saída em (T-If) é \(\Gamma\), e \textbf{não}
\(\Gamma_1\) nem \(\Gamma_2\). Os contextos \(\Gamma_1\) e \(\Gamma_2\)
são descartados ao sair dos blocos, de modo que variáveis locais a
cada ramo não escapam para o escopo externo. O mesmo vale para
(T-While): \(\Gamma'\) é descartado e o contexto retorna a \(\Gamma\).

% ============================================================
\section{Sistema de Tipos para Funções e Registros}
% ============================================================

A linguagem TImp estende TWhile com declarações de funções e de tipos
registro. Essas duas extensões exigem enriquecer os ambientes de
tipagem e adicionar novas regras ao sistema. Nesta seção
apresentamos o sistema de tipos de TImp de forma completa.

% ------------------------------------------------------------
\subsection{Tipos, Contextos e Ambientes}
% ------------------------------------------------------------

\paragraph{Tipos.}

\[
  \tau \;::=\; \mathbf{Int} \mid \mathbf{Bool} \mid \mathbf{String}
               \mid R
  \qquad\qquad
  \rho \;::=\; \mathbf{void} \mid \tau
\]

O tipo \(R\) é um \textbf{tipo registro} nomeado; \(\rho\) é o tipo de
retorno de funções: ou o tipo unitário \(\mathbf{void}\) (sem valor
de retorno) ou um tipo \(\tau\).

\paragraph{Contexto de variáveis.}

\[
  \Gamma : \mathit{Var} \rightharpoonup \tau
\]

\paragraph{Ambiente de registros.}

\[
  \Delta : \mathit{Name} \rightharpoonup \overline{(f_i : \tau_i)}
\]

\(\Delta(R)\) lista os campos e seus tipos para o tipo registro \(R\).
Escreve-se \(\Delta(R)(f)\) para o tipo do campo \(f\) em \(R\).

\paragraph{Ambiente de funções.}

\[
  \Phi : \mathit{Name} \rightharpoonup ([\tau_1,\ldots,\tau_n],\; \rho)
\]

\(\Phi(f)\) fornece a lista de tipos dos parâmetros e o tipo de
retorno de \(f\). Ambos \(\Delta\) e \(\Phi\) são construídos das
declarações do programa e permanecem \emph{imutáveis} durante a
verificação.

A relação de tipagem de expressões torna-se
\(\Gamma;\Delta;\Phi \vdash e : \tau\).
A de comandos torna-se \(\Gamma;\Delta;\Phi;\rho_{\mathit{ret}}
\vdash s \dashv \Gamma'\), onde \(\rho_{\mathit{ret}}\) é o tipo de
retorno esperado da função corrente (ou \(\mathbf{void}\) no nível
superior). Por brevidade, omitimos \(\Delta\), \(\Phi\) e
\(\rho_{\mathit{ret}}\) quando não são necessários numa regra.

% ------------------------------------------------------------
\subsection{Regras de Tipagem para Registros}
% ------------------------------------------------------------

\subsubsection{Expressões}

\textbf{Acesso a campo:}

\[
  \infer[\text{(T-Field)}]
  {\Gamma;\Delta;\Phi \vdash e.f : \tau}
  {\Gamma;\Delta;\Phi \vdash e : R
   \quad R \in \mathrm{dom}(\Delta)
   \quad \Delta(R)(f) = \tau}
\]

\textbf{Construção de registro:}

\[
  \infer[\text{(T-New)}]
  {\Gamma;\Delta;\Phi \vdash \mathbf{new}\; R\,\{f_1\!=\!e_1,\ldots,f_n\!=\!e_n\} : R}
  {R \in \mathrm{dom}(\Delta)
   \quad \Delta(R) = \overline{(f_i:\tau_i)}
   \quad \Gamma;\Delta;\Phi \vdash e_i : \tau_i \;\;(\forall\,i)}
\]

A regra (T-New) exige que o conjunto de campos fornecidos coincida
exatamente com o conjunto de campos declarados em \(\Delta(R)\) e que
cada campo tenha o tipo correto.

\subsubsection{Comandos}

\textbf{Atribuição de campo:}

\[
  \infer[\text{(T-FieldAssign)}]
  {\Gamma;\Delta;\Phi \vdash x.f \mathbin{:=} e \dashv \Gamma}
  {\Gamma(x) = R
   \quad \Delta(R)(f) = \tau
   \quad \Gamma;\Delta;\Phi \vdash e : \tau}
\]

% ------------------------------------------------------------
\subsection{Regras de Tipagem para Funções}
% ------------------------------------------------------------

\subsubsection{Chamadas de função}

\textbf{Chamada como expressão} (função não-\texttt{void}):

\[
  \infer[\text{(T-CallExpr)}]
  {\Gamma;\Delta;\Phi \vdash f(e_1,\ldots,e_n) : \tau}
  {\Phi(f) = ([\tau_1,\ldots,\tau_n],\; \tau)
   \quad \Gamma;\Delta;\Phi \vdash e_i : \tau_i \;\;(\forall\,i)}
\]

\textbf{Chamada como comando} (função \texttt{void}):

\[
  \infer[\text{(T-CallStmt)}]
  {\Gamma;\Delta;\Phi \vdash f(e_1,\ldots,e_n) \dashv \Gamma}
  {\Phi(f) = ([\tau_1,\ldots,\tau_n],\; \mathbf{void})
   \quad \Gamma;\Delta;\Phi \vdash e_i : \tau_i \;\;(\forall\,i)}
\]

A distinção entre (T-CallExpr) e (T-CallStmt) é essencial: só funções
\texttt{void} podem ser usadas como comandos; usar uma função
não-\texttt{void} como comando desperdiçaria o valor de retorno e é
considerado um erro de tipo.

\subsubsection{Retorno}

As regras de tipagem para \textbf{return} dependem do tipo de retorno
\(\rho_{\mathit{ret}}\) da função corrente, propagado pelo ambiente de
tipagem:

\[
  \infer[\text{(T-RetVal)}]
  {\Gamma;\ldots;\rho_{\mathit{ret}} \vdash \mathbf{return}\; e \dashv \Gamma}
  {\rho_{\mathit{ret}} = \tau
   \quad \Gamma;\Delta;\Phi \vdash e : \tau}
\]

\[
  \infer[\text{(T-RetVoid)}]
  {\Gamma;\ldots;\rho_{\mathit{ret}} \vdash \mathbf{return} \dashv \Gamma}
  {\rho_{\mathit{ret}} = \mathbf{void}}
\]

\subsubsection{Declaração de função}

A tipagem de uma declaração de função verifica o corpo sob um contexto
inicial formado pelos parâmetros formais, com \(\rho_{\mathit{ret}}\)
fixado ao tipo de retorno declarado:

\[
  \infer[\text{(T-FuncDecl)}]
  {\Delta;\Phi \vdash \mathbf{fn}\; f(\overline{x_i:\tau_i}): \rho\; \{B\} \;\checkmark}
  {[x_1\!\mapsto\!\tau_1,\ldots,x_n\!\mapsto\!\tau_n];\Delta;\Phi;\rho
   \vdash B \dashv \_}
\]

O contexto de saída do bloco é descartado: parâmetros e variáveis
locais à função não são visíveis no nível superior.

% ------------------------------------------------------------
\subsection{Exemplo de Derivação de Tipo}
% ------------------------------------------------------------

Considere a expressão \(\mathbf{new}\;
\mathtt{Point}\,\{x\!=\!3,\, y\!=\!4\}.x + 1\) com
\(\Delta(\mathtt{Point}) = [(x:\mathbf{Int}),\, (y:\mathbf{Int})]\) e
\(\Gamma = \emptyset\):

\[
  \dfrac{
    \dfrac{
      \dfrac{
        \dfrac{}{\vdash 3 : \mathbf{Int}}
        \quad
        \dfrac{}{\vdash 4 : \mathbf{Int}}
      }{\vdash \mathbf{new}\;\mathtt{Point}\{x\!=\!3,y\!=\!4\} : \mathtt{Point}}
      \quad \Delta(\mathtt{Point})(x) = \mathbf{Int}
    }{\vdash (\mathbf{new}\;\mathtt{Point}\{x\!=\!3,y\!=\!4\}).x : \mathbf{Int}}
    \quad
    \dfrac{}{\vdash 1 : \mathbf{Int}}
  }{\vdash (\mathbf{new}\;\mathtt{Point}\{x\!=\!3,y\!=\!4\}).x + 1 : \mathbf{Int}}
\]

% ------------------------------------------------------------
\subsection{Extensão da Corretude}
% ------------------------------------------------------------

As propriedades de preservação e de progresso estendem-se
naturalmente para TImp. Dois aspectos merecem nota:

\begin{itemize}
  \item \textbf{Preservação} para (T-Field) requer que o valor
    \(\mathbf{rec}(R, \mathit{flds})\) tenha os campos esperados por
    \(\Delta(R)\), invariante mantido por (T-New) e (T-FieldAssign).
  \item \textbf{Progresso} para chamadas de função requer que o corpo
    de \(f\) termine com um \(\mathbf{return}\; e\) quando \(\rho \neq
    \mathbf{void}\). A regra (T-FuncDecl) tipifica o corpo, mas a
    garantia de que toda execução efetivamente atinge um
    \haskell{return} exige \textbf{análise de fluxo de controlo}
    (verificação de retorno garantido), que está além do escopo deste
    capítulo.
\end{itemize}

% ============================================================
\section{Conclusão}

Neste capítulo, estudamos os fundamentos dos sistemas de tipos em
linguagens de programação. Vimos que tipos classificam os valores de
uma linguagem e restringem as operações admissíveis, e que a
verificação de tipos pode ocorrer em tempo de compilação (tipagem
estática) ou em tempo de execução (tipagem dinâmica).
Ortogonalmente, uma linguagem pode ser fortemente tipada, proibindo
conversões implícitas entre tipos incompatíveis, ou fracamente
tipada, admitindo tais coerções de forma automática.

Para dar uma base formal a esses conceitos, utilizamos a linguagem de
expressões tipadas de Pierce, que combina booleanos e números
naturais. O seu sistema de tipos é definido por regras de inferência
da forma \(\vdash t : T\), que atribuem um tipo a cada termo de
maneira única, propriedade expressa pelo lema de unicidade de
tipos. A semântica small-step descreve a execução como uma sequência
de reduções \(t \to t'\) e torna explícito o conceito de
\textbf{formas presas} (\emph{stuck terms}): termos que não são
valores e para os quais nenhuma regra de redução se aplica,
representando erros irrecuperáveis em tempo de execução. A semântica
big-step, por sua vez, relaciona cada termo diretamente ao seu valor
final \(t \Downarrow v\), abstraindo os passos intermediários.

A propriedade central que justifica a utilidade da tipagem estática é
a \textbf{corretude do sistema de tipos} (\emph{type soundness}):
programas bem tipados nunca ficam presos. Demonstramos essa
propriedade pela composição de dois lemas independentes. O lema de
\textbf{preservação} garante que cada passo de redução mantém o tipo
do termo: a informação de tipo não se corrompe durante a execução.
O lema de \textbf{progresso} garante que um termo bem tipado que
ainda não é valor sempre pode dar um passo de redução, de modo que a execução
nunca trava. Juntos, os dois lemas asseguram que toda computação a
partir de um termo bem tipado avança indefinidamente ou termina em um
valor, nunca numa forma presa.

Estendemos o formalismo para linguagens imperativas apresentando os
sistemas de tipos de TLine, TWhile e TImp. Para TLine, a presença de
variáveis exige um contexto \(\Gamma\) e as regras distinguem
declaração (que estende \(\Gamma\)) de atribuição (que o preserva).
Para TWhile, a adição de condicionais e repetição com escopo léxico é
formalizada pelas regras (T-If) e (T-While), que mantêm o contexto de
saída igual ao de entrada, descartando variáveis declaradas internamente.
Para TImp, funções e registros introduzem dois ambientes globais
imutáveis (\(\Delta\) para registros e \(\Phi\) para funções),
além do tipo de retorno esperado \(\rho_{\mathit{ret}}\) por função. A
implementação desses verificadores de tipos é tratada no próximo capítulo.
